% ==================================================================
\section{Recursive: Improved Leaf Task Decomposition}
\label{sec:rec-improved}
% ==================================================================

\subsection{Implementation}
\label{sec:rec-improved-impl}

The \textbf{improved leaf} version addresses the main bottleneck of the plain leaf
strategy: serial task generation.
In the leaf version, the entire recursion is driven by a single thread, so threads
sit idle until enough leaf tasks have accumulated.
The improvement introduces a \emph{depth cutoff} parameter \texttt{d} (initialised to
\texttt{0} at the top-level call) and a compile-time constant \texttt{CUTOFF}.

\begin{itemize}
    \item When the current recursion depth satisfies $\mathtt{d} < \mathtt{CUTOFF}$,
          each recursive quadrant call is wrapped in a \texttt{\#pragma~omp~task},
          allowing multiple threads to explore different branches of the recursion tree
          \emph{concurrently}.
    \item When $\mathtt{d} \ge \mathtt{CUTOFF}$, the recursive calls are sequential
          again, cutting off task proliferation at deep levels where the overhead would
          dominate.
    \item Leaf tasks (fill or compute) are only spawned when $\mathtt{d} \le
          \mathtt{CUTOFF}$; beyond the cutoff they run inline to avoid the overhead of
          tiny tasks.
\end{itemize}

The value \texttt{CUTOFF = 16} was chosen empirically so that the number of tasks at
the cutoff level matches or exceeds the number of available threads.

\begin{sloppypar}
The complete implementation (depth-cutoff mechanism with \texttt{CUTOFF = 16}) is in \path{codis/leaf-improved.c}.
\end{sloppypar}

The \texttt{\#pragma~omp~taskgroup} ensures that all four quadrant tasks (and their
descendants) finish before control returns to the caller, preserving the border-read
correctness that the leaf version relies on.

\paragraph{Improvements over plain Leaf.}
\begin{itemize}
    \item Multiple threads now explore the recursion tree in parallel down to depth
          \texttt{CUTOFF}, eliminating the serial task-generation bottleneck.
    \item Tasks beyond \texttt{CUTOFF} run sequentially, avoiding excessive fine-grained
          task overhead at the bottom of the tree.
    \item The result is speedups larger than $2\times$ for more than 2 threads.
\end{itemize}

\subsection{Modelfactor Analysis}
\label{sec:rec-improved-mf}

Tables~\ref{tab:li-exec}--\ref{tab:li-tasks} summarise the Modelfactors output for
\texttt{mandel-omp-rec-improved-leaf} obtained on \texttt{par2317}.

\begin{table}[h]
\resizebox{\textwidth}{!}{%
\begin{tabular}{|l|c|c|c|c|c|c|c|c|c|c|}
\hline
\multicolumn{11}{|c|}{Overview of whole program execution metrics} \\
\hline
\hline
Number of threads & 1 & 2 & 4 & 8 & 12 & 16 & 20 & 24 & 28 & 32 \\
\hline
Elapsed time (sec)      &       1.82 &       0.99 &       0.58 &       0.38 &       0.38 &       0.38 &       0.38 &       0.38 &       0.38 &       0.38 \\
\hline
Speedup                 &       1.00 &       1.84 &       3.12 &       4.76 &       4.76 &       4.74 &       4.74 &       4.74 &       4.73 &       4.74 \\
\hline
Efficiency              &       1.00 &       0.92 &       0.78 &       0.60 &       0.40 &       0.30 &       0.24 &       0.20 &       0.17 &       0.15 \\
\hline
\end{tabular}
}
\caption{Improved Leaf -- whole program execution metrics. Analysis: par2317, 2026-05-15.}
\label{tab:li-exec}
\end{table}


\begin{table}[h]
\resizebox{\textwidth}{!}{%
\begin{tabular}{|l|c|c|c|c|c|c|c|c|c|c|}
\hline
\multicolumn{11}{|c|}{Overview of the Efficiency metrics in parallel fraction, $\phi$=99.98\%} \\
\hline
\hline
Number of threads & 1 & 2 & 4 & 8 & 12 & 16 & 20 & 24 & 28 & 32 \\
\hline
\hline
Global efficiency                      &     99.96\% &     91.86\% &     77.89\% &     59.57\% &     39.63\% &     29.65\% &     23.69\% &     19.77\% &     16.91\% &     14.81\% \\
\hline
\hline
Parallelization strategy efficiency &     99.96\% &     92.02\% &     78.01\% &     59.70\% &     39.78\% &     29.83\% &     23.91\% &     20.00\% &     17.18\% &     15.11\% \\
\hline
Load balancing                   &    100.00\% &     92.07\% &     80.14\% &     66.58\% &     45.42\% &     31.87\% &     27.40\% &     22.91\% &     18.37\% &     17.31\% \\
In execution efficiency          &     99.96\% &     99.94\% &     97.34\% &     89.66\% &     87.57\% &     93.59\% &     87.25\% &     87.31\% &     93.54\% &     87.31\% \\
\hline
\hline
Scalability for computation tasks   &    100.00\% &     99.83\% &     99.85\% &     99.78\% &     99.63\% &     99.40\% &     99.10\% &     98.82\% &     98.42\% &     98.04\% \\
\hline
IPC scalability                  &    100.00\% &     99.94\% &     99.95\% &     99.97\% &     99.94\% &     99.92\% &     99.90\% &     99.93\% &     99.93\% &     99.91\% \\
Instruction scalability          &    100.00\% &    100.00\% &    100.00\% &    100.00\% &    100.00\% &    100.00\% &    100.00\% &    100.00\% &    100.00\% &    100.00\% \\
Frequency scalability            &    100.00\% &     99.89\% &     99.89\% &     99.82\% &     99.68\% &     99.48\% &     99.19\% &     98.89\% &     98.49\% &     98.12\% \\
\hline
\end{tabular}
}
\caption{Improved Leaf -- efficiency breakdown ($\phi$=99.98\%). Analysis: par2317, 2026-05-15.}
\label{tab:li-efficiency}
\end{table}


\begin{table}[h]
\resizebox{\textwidth}{!}{%
\begin{tabular}{|l|c|c|c|c|c|c|c|c|c|c|}
\hline
\multicolumn{11}{|c|}{Statistics about explicit tasks in parallel fraction} \\
\hline
\hline
Number of threads & 1 & 2 & 4 & 8 & 12 & 16 & 20 & 24 & 28 & 32 \\
\hline
\hline
Number of explicit tasks executed (total)        &            66.0 &            66.0 &            66.0 &            66.0 &            66.0 &            66.0 &            66.0 &            66.0 &            66.0 &            66.0 \\
\hline
LB (number of explicit tasks executed)           &             1.0 &            0.87 &            0.69 &            0.49 &            0.55 &            0.59 &            0.55 &            0.55 &            0.47 &            0.55 \\
\hline
LB (time executing explicit tasks)               &             1.0 &            0.93 &            0.79 &            0.69 &            0.55 &            0.47 &             0.4 &             0.3 &            0.24 &            0.27 \\
\hline
Time per explicit task (average us)                 &        27349.36 &        27384.63 &        27386.04 &        28433.09 &        34192.49 &        40782.88 &        41755.11 &        37641.56 &        37773.34 &        41457.39 \\
\hline
Overhead per explicit task (synch \%)             &            0.04 &            8.71 &           28.37 &           65.43 &          122.21 &          159.14 &          210.86 &          295.23 &          355.38 &          379.16 \\
\hline
Overhead per explicit task (sched \%)             &             0.0 &             0.0 &            0.01 &            0.01 &            0.02 &            0.01 &            0.02 &            0.01 &            0.02 &            0.02 \\
\hline
Number of taskwait/taskgroup (total)             &          1004.0 &          1004.0 &          1004.0 &          1004.0 &          1004.0 &          1004.0 &          1004.0 &          1004.0 &          1004.0 &          1004.0 \\
\hline
\end{tabular}
}
\caption{Improved Leaf -- explicit task statistics. Analysis: par2317, 2026-05-15.}
\label{tab:li-tasks}
\end{table}


\paragraph{Observations.}
The improved leaf version exposes a clear performance ceiling.
Elapsed time decreases from 1.82\,s (1 thread) to 0.38\,s at 8 threads
(speedup $4.76\times$) and then remains \emph{flat} all the way to 32 threads
(Table~\ref{tab:li-exec}).
Efficiency collapses to 15\,\% at 32 threads.

The efficiency breakdown (Table~\ref{tab:li-efficiency}) identifies \emph{load
imbalance} as the dominant bottleneck: the load-balancing factor falls from 100\,\% at
1 thread to only 17.31\,\% at 32 threads.
This is a direct consequence of the small total task count: only \textbf{66 explicit
tasks} are generated (Table~\ref{tab:li-tasks}), making it impossible to distribute
work evenly across more than a handful of threads.

The synchronisation overhead per task tells the same story:
it grows from 0.04\,\% at 1 thread to \textbf{379\,\%} at 32 threads, meaning threads
spend nearly four times as long waiting at synchronisation points (\texttt{taskwait}/
\texttt{taskgroup}) as they do executing useful computation.
This is caused by the 1\,004 \texttt{taskwait}/\texttt{taskgroup} operations that must
serialise access between recursion levels, each becoming a bottleneck when more threads
contend for a small pool of tasks.
Scheduling overhead, by contrast, is negligible ($\leq$0.02\,\%), confirming that the
overhead is synchronisation-bound, not dispatch-bound.

\begin{sloppypar}
The saturation at 8 threads matches the CUTOFF-induced task count: with \texttt{CUTOFF=16},
the recursion tree generates at most $4^{\lceil\log_4 66\rceil}$ leaf tasks, and at 8
threads the work is fully distributed; additional threads find no tasks and remain idle.
\end{sloppypar}

Compared with the plain Leaf version, the improved leaf eliminates the serial
task-generation phase (multiple threads now traverse the recursion tree in parallel up
to the cutoff), but the fundamental limit of 66 tasks prevents further scaling beyond
8 threads.
Increasing \texttt{CUTOFF} would generate more tasks and potentially improve scalability
at the cost of higher synchronisation overhead per task.

\subsection{Paraver Analysis}
\label{sec:rec-improved-paraver}

% TODO: Include Paraver views for 16 threads.
%   Compare with the plain Leaf trace:
%   - Is the initial serial phase shorter or eliminated?
%   - Is load balance better?

\todo[inline]{Insert Paraver screenshots (16 threads) for Improved Leaf.
Include both the non-improved and the improved traces for direct comparison.}

% \begin{figure}[H]
%   \centering
%   \begin{subfigure}[b]{0.48\linewidth}
%     \includegraphics[width=\linewidth]{fotos/rec-leaf-paraver-parallelism.png}
%     \caption{Plain Leaf -- parallelism view.}
%   \end{subfigure}
%   \hfill
%   \begin{subfigure}[b]{0.48\linewidth}
%     \includegraphics[width=\linewidth]{fotos/rec-improved-paraver-parallelism.png}
%     \caption{Improved Leaf -- parallelism view.}
%   \end{subfigure}
%   \caption{Paraver parallelism comparison, 16 threads.}
%   \label{fig:improved-paraver}
% \end{figure}

\subsection{Strong Scalability}
\label{sec:rec-improved-scalability}

% TODO: Include speedup plot for improved leaf, compared with plain leaf.

\todo[inline]{Insert the strong scalability plot for Improved Leaf and compare with
plain Leaf.  Confirm that speedups larger than $2\times$ are achieved for more than
2 threads.}

% \begin{figure}[H]
%   \centering
%   \includegraphics[width=0.75\linewidth]{fotos/rec-improved-speedup.png}
%   \caption{Strong scalability: Recursive Leaf vs.\ Improved Leaf.}
%   \label{fig:improved-speedup}
% \end{figure}
